% ========================================================
% ANÁLISIS DEL CÓDIGO CON IA: GAPs (ENTREGABLE 3)
% ========================================================
% El informe completo es un .md, como pide el enunciado: tp/docs/gaps-ia.md. Esta sección lo
% resume y explica el método; el detalle técnico queda en el .md (Anexo B).
\section{Análisis del código con IA: GAPs}\label{sec:gaps}

El enunciado pide elaborar un prompt estructurado, elegir un modelo de IA, analizar con él el código
del repositorio y documentar en Markdown los GAPs en calidad de código, seguridad, patrones de
concurrencia y otros aspectos relevantes. El informe completo está en
\texttt{tp/docs/gaps-ia.md} (Anexo~\ref{anexo:gaps}); esta sección resume el método y los hallazgos.

\subsection{Modelo y prompt}

\pendiente{qué modelo se eligió, con qué versión y por qué (contexto largo para leer el repositorio,
capacidad para razonar sobre concurrencia, disponibilidad); con qué herramienta se le dio acceso al
código y a qué commit}

\pendiente{la estructura del prompt: rol, contexto del proyecto, alcance (qué directorios), criterios
por categoría, formato de salida exigido (hallazgo, evidencia con archivo y línea, severidad,
recomendación) y la instrucción de no inventar hallazgos sin evidencia. El prompt completo va en el
.md}

\subsection{Hallazgos}

\pendiente{tabla resumen: categoría (calidad de código, seguridad, patrones de concurrencia,
pruebas, rendimiento, documentación), cantidad de hallazgos y el más importante de cada una}

\subsection{Contraste con la verificación del proyecto}

\pendiente{qué hallazgos de la IA se confirmaron contra el código y cuáles no (falsos positivos),
qué ya estaba cubierto por \texttt{go test -race} o por Spin, y qué encontró la IA que ninguna de las
dos herramientas cubre. El sílabo pide contrastar lo que dice la IA con la teoría: esta subsección es
ese contraste}
